% Part: lambda-calculus
% Chapter: introduction
% Section: lambda-definability

\documentclass[../../../include/open-logic-section]{subfiles}

\begin{document}

\olfileid[id]{lam}{int}{rep}
\olsection{Fungsi Aritmetika yang \usetoken{S}{lambda definable}}

Bagaimana kalkulus lambda dapat berfungsi sebagai model komputasi? Pada
awalnya, bahkan belum jelas bagaimana pernyataan ini harus dipahami. Untuk
membahas komputabilitas pada bilangan asli, kita perlu menemukan representasi
yang sesuai bagi bilangan-bilangan tersebut. Berikut adalah representasi yang
ternyata bekerja dengan sangat baik.

\begin{defn}
Untuk setiap bilangan asli~$n$, definisikan \emph{numeral Church} $\num{n}$
sebagai suku lambda $\lambd[x][\lambd[y][(x(x(x(\dots x(y)))))]]$, dengan
tepat $n$ kemunculan~$x$.
\end{defn}

Suku-suku $\num{n}$ merupakan ``iterator'': pada masukan~$f$, $\num{n}$
mengembalikan fungsi yang memetakan $y$ ke $f^n(y)$. Perhatikan bahwa setiap
numeral berbentuk normal. Sekarang kita dapat menyatakan apa artinya suatu suku
lambda ``menghitung'' fungsi pada bilangan asli.

\begin{defn}
Misalkan $f(x_0, \dots, x_{k-1})$ adalah fungsi parsial beraritas~$k$ dari
$\Nat^k$ ke~$\Nat$. Kita mengatakan bahwa suatu suku-$\lambd$~$F$
\emph{!!{lambda define}s}~$f$ jika dan hanya jika, untuk setiap barisan bilangan
asli $n_0$, \dots,~$n_{k-1}$,
\[
F\, \num{n_0}\, \num{n_1} \dots \num{n_{k-1}} \red \num{f(n_0, n_1, \dots,
  n_{k-1})}
\]
apabila $f(n_0, \dots, n_{k-1})$ terdefinisi, sedangkan
$F\,\num{n_0}\,\num{n_1}\dots\num{n_{k-1}}$ tidak memiliki bentuk normal
dalam kasus lainnya.
\end{defn}

\begin{thm}
\ollabel{thm:lambda-def}
Suatu fungsi $f$ merupakan fungsi dapat dihitung parsial jika dan hanya jika
fungsi itu !!{lambda defined} oleh suatu suku lambda.
\end{thm}

\begin{explain}
Teorema ini cukup mencolok. Sebagai model komputasi, kalkulus lambda merupakan
kalkulus yang sangat sederhana; satu-satunya operasinya adalah abstraksi lambda
dan aplikasi!{} Namun, dengan sumber daya yang terbatas ini kita tetap dapat
mengimplementasikan setiap prosedur komputasi.
\end{explain}

\end{document}
